
\documentclass[12pt,a4paper]{article}

\usepackage[a4paper,margin=1in]{geometry}
\usepackage{graphicx}
\usepackage{booktabs}
\usepackage{amsmath,amssymb}
\usepackage{float}
\usepackage{hyperref}
\usepackage{xcolor}
\usepackage{listings}
\usepackage{enumitem}
\usepackage{fancyhdr}
\usepackage{longtable}

\hypersetup{colorlinks=true,linkcolor=blue,urlcolor=blue}

\pagestyle{fancy}
\fancyhf{}
\lhead{ICS1512}
\rhead{Machine Learning Algorithms Laboratory}
\cfoot{\thepage}

\lstset{
language=Python,
basicstyle=\ttfamily\small,
numbers=left,
frame=single,
breaklines=true,
showstringspaces=false
}

\begin{document}

\begin{tabular}{|p{4cm}|p{11cm}|}
\hline
Experiment No. & 8 \\ \hline
Experiment Title & Clustering Human Activity Recognition Data using K-Means, DBSCAN, and Hierarchical
Clustering\footnotemark \\ \hline
Student Name & Sanjay S \\ \hline
Register Number & 3122247001056 \\ \hline
Faculty & Dr. Poreddy Ajay Kumar Reddy \\ \hline
Submission Date & 21/09/2026 \\ \hline
\end{tabular}
\footnotetext{\textbf{GitHub Repository: }\url{https://github.com/sanjay-6727/Introdution-to-Machine-Learning-}}

\section{Objective}
Try clustering on the Human Activity Recognition dataset three different ways -- K-Means, DBSCAN and
Hierarchical Agglomerative Clustering -- and actually see how each one carves up the data compared to
the six real activity labels, instead of just assuming clustering will neatly rediscover the activities
on its own.

\section{Problem Statement}
Given the UCI HAR smartphone sensor dataset, preprocess and standardize it, reduce it for
visualization, then run K-Means (choosing $k$ with the elbow method), DBSCAN (tuning $\epsilon$ and
minPts), and Agglomerative Clustering (comparing linkage methods) on it. Compare the three using
internal metrics (silhouette, Davies-Bouldin, Calinski-Harabasz) and external metrics against the true
activity labels (ARI, NMI), and figure out which algorithm actually produces the most meaningful
clusters here.

\section{Dataset Description}
\begin{longtable}{|p{3.5cm}|p{4cm}|p{1.6cm}|p{1.8cm}|p{2cm}|p{2.2cm}|}
\hline
\textbf{Dataset} & \textbf{Source} & \textbf{Samples} & \textbf{Features} & \textbf{Activities} & \textbf{Subjects} \\ \hline
Human Activity Recognition Using Smartphones & UCI Machine Learning Repository & 10299 (7352 train +
2947 test) & 561 & 6 & 30 \\ \hline
\end{longtable}

\textbf{Activity distribution:} LAYING = 1944, STANDING = 1906, SITTING = 1777, WALKING = 1722,
WALKING\_UPSTAIRS = 1544, WALKING\_DOWNSTAIRS = 1406. Reasonably balanced across the six classes, with
LAYING slightly over-represented and WALKING\_DOWNSTAIRS slightly under.

\textbf{Preprocessing steps:}
\begin{itemize}
\item Merged the official train and test splits back into one 10299-row dataset, since clustering is
unsupervised and doesn't need the original train/test split, and combining the two just gives a fuller
picture of the activity space.
\item No missing values anywhere in the 561 features, so no imputation was needed.
\item Activity codes (1--6) mapped to their proper names using \texttt{activity\_labels.txt}, kept
around purely as ground truth for the external evaluation metrics, never fed to the clustering
algorithms themselves.
\item All 561 features standardized with \texttt{StandardScaler}, since clustering algorithms are
distance-based and the raw features already sit in roughly $[-1, 1]$ but on slightly different scales.
\item Since 10299 points is a lot to run DBSCAN, Agglomerative clustering and t-SNE on comfortably, I
took a stratified 2000-point subsample (proportional per activity) for the actual clustering and
visualization, so every algorithm sees the same data and the comparison stays fair.
\item PCA down to 30 components was used as the working feature space for all three clustering
algorithms (about 82\% of the variance retained), mainly to keep distances meaningful and DBSCAN/HAC
tractable; a separate 2-component PCA and a 2D t-SNE embedding were computed purely for plotting.
\end{itemize}

\section{Model A: K-Means Clustering and the Elbow Method}

\begin{table}[H]
\centering
\caption{K-Means Elbow Method Results}
\begin{tabular}{ccc}
\toprule
Number of Clusters (k) & WCSS (Inertia) & Silhouette Score \\
\midrule
2 & 433577.6 & 0.4775 \\
3 & 359760.2 & 0.4070 \\
4 & 332567.7 & 0.2174 \\
5 & 305593.0 & 0.1940 \\
6 & 291478.2 & 0.1835 \\
7 & 280425.3 & 0.1602 \\
8 & 270008.5 & 0.1573 \\
\bottomrule
\end{tabular}
\end{table}

\begin{figure}[H]
\centering
\includegraphics[width=\linewidth]{outputs/elbow_silhouette.png}
\caption{Elbow curve (WCSS vs k) and silhouette score vs k for K-Means.}
\end{figure}

Purely on silhouette score, $k=2$ wins by a wide margin, and the elbow curve also bends hardest around
$k=2$ to $k=3$. That's not surprising -- at the coarsest level, HAR activities split into two very
distinct physical states, moving (the three walking types) and stationary (sitting, standing, laying),
and geometrically those two blobs are far apart while the finer distinctions inside each blob are much
subtler. Since the assignment goal is to compare clustering against the six known activities, though,
I used \textbf{k=6} for the main analysis and all cross-algorithm comparisons, so K-Means, DBSCAN and
Hierarchical Clustering are all being judged on the same basis.

\section{Model B: DBSCAN}

minPts was fixed at 60 (twice the 30-dimensional PCA space, a common rule of thumb), and $\epsilon$ was
swept using a k-distance plot plus a direct search over several candidate values.

\begin{figure}[H]
\centering
\includegraphics[width=0.65\linewidth]{outputs/dbscan_kdistance.png}
\caption{60-NN distance plot used to guide the choice of $\epsilon$ for DBSCAN.}
\end{figure}

\begin{table}[H]
\centering
\caption{DBSCAN -- eps Tuning}
\begin{tabular}{cccc}
\toprule
eps & Clusters Found & Noise Points & Silhouette (non-noise) \\
\midrule
8 & 1 & 1220 (61.0\%) & -- \\
10 & 1 & 996 (49.8\%) & -- \\
12 & 2 & 483 (24.1\%) & 0.5011 \\
14 & 1 & 219 (11.0\%) & -- \\
16 & 1 & 108 (5.4\%) & -- \\
18 & 1 & 70 (3.5\%) & -- \\
20 & 1 & 37 (1.9\%) & -- \\
\bottomrule
\end{tabular}
\end{table}

$\epsilon = 12$ was the only setting in the sweep that actually produced more than one cluster; every
other value either merged everything into one blob (too permissive) or scattered most points as noise
(too strict). The final DBSCAN run at $\epsilon = 12$, minPts $=60$ gives \textbf{2 clusters} and 483
noise points (24.1\% of the data).

\section{Model C: Hierarchical Agglomerative Clustering}

\begin{table}[H]
\centering
\caption{Hierarchical Clustering -- Linkage Comparison (k=6)}
\begin{tabular}{ccc}
\toprule
Linkage & Silhouette Score & Davies-Bouldin Index \\
\midrule
Single & 0.5590 & 0.4092 \\
Complete & 0.4458 & 1.0389 \\
Average & 0.4625 & 0.9031 \\
Ward & 0.1663 & 1.8512 \\
\bottomrule
\end{tabular}
\end{table}

\begin{figure}[H]
\centering
\includegraphics[width=\linewidth]{outputs/dendrogram_ward.png}
\caption{Dendrogram using Ward's linkage on a 200-point subsample.}
\end{figure}

Single linkage posts the best silhouette score here, but that's a classic chaining artifact -- it
tends to swallow almost everything into one giant cluster and peel off a handful of outliers as tiny
separate clusters, which looks great by silhouette but doesn't map onto anything meaningful. Ward's
linkage, which the experiment specifically asks for, gives a much lower silhouette (0.1663) but that's
because it's actually trying to build six balanced, compact clusters instead of one big blob plus
scraps, so it's the more honest and more useful result even though the raw silhouette number looks
worse on paper. Ward's linkage is what's used for the rest of the comparison below.

\section{Evaluation Metrics and Results}

\begin{table}[H]
\centering
\caption{Clustering Evaluation Metrics Comparison}
\begin{tabular}{lccccc}
\toprule
Model & Silhouette & Davies-Bouldin & Calinski-Harabasz & ARI & NMI \\
\midrule
K-Means (k=6) & 0.1835 & 1.7524 & 858.3 & 0.4178 & 0.5629 \\
DBSCAN (eps=12) & 0.5011 & 0.8295 & 1810.5 & 0.3022 & 0.5321 \\
Hierarchical (Ward) & 0.1663 & 1.8512 & 794.6 & 0.3541 & 0.5345 \\
\bottomrule
\end{tabular}
\end{table}

\begin{figure}[H]
\centering
\includegraphics[width=0.75\linewidth]{outputs/metrics_barplot.png}
\caption{Silhouette, ARI and NMI compared across the three clustering algorithms.}
\end{figure}

K-Means comes out ahead on the metrics that actually compare against the true activity labels (ARI
0.4178, NMI 0.5629), even though DBSCAN wins on the purely internal silhouette and Davies-Bouldin
scores. That's the DBSCAN result being a bit of a mirage: it only found 2 clusters (moving vs
stationary), and it's easy to get a great internal separation score with just 2 well-separated blobs
while completely missing the finer 6-way structure the experiment actually cares about.

\section{Graphs and Visualization}

\begin{figure}[H]
\centering
\includegraphics[width=\linewidth]{outputs/ground_truth_scatter.png}
\caption{PCA (left) and t-SNE (right) 2D projections colored by the true activity label, for
reference.}
\end{figure}

\begin{figure}[H]
\centering
\includegraphics[width=\linewidth]{outputs/cluster_scatter_pca_tsne.png}
\caption{PCA (top row) and t-SNE (bottom row) projections colored by cluster assignment for K-Means,
DBSCAN and Hierarchical (Ward) respectively.}
\end{figure}

\begin{figure}[H]
\centering
\includegraphics[width=0.75\linewidth]{outputs/kmeans_vs_truth_confusion.png}
\caption{K-Means clusters (k=6) cross-tabulated against the true activity labels.}
\end{figure}

\textbf{Inference:} The ground-truth t-SNE plot shows a very clean, tight LAYING blob sitting far from
everything else, a fairly clean WALKING/WALKING\_UPSTAIRS/WALKING\_DOWNSTAIRS region, and SITTING and
STANDING overlapping quite a bit with each other. That overlap between SITTING and STANDING is exactly
where all three clustering algorithms struggle: in the K-Means cross-tab, cluster 2 mixes 244 SITTING
with 249 STANDING points almost 50/50, and cluster 5 mixes 88 SITTING, 121 STANDING and 56 LAYING
together. LAYING, by contrast, gets almost cleanly isolated into cluster 0 (305 of 378 LAYING points),
and the three walking activities mostly separate out into clusters 1, 3 and 4, though
WALKING\_UPSTAIRS and WALKING\_DOWNSTAIRS bleed into each other somewhat too.

\section{Observations}
\begin{itemize}
\item \textbf{Which algorithm produced the most meaningful clusters:} K-Means, on balance. It's the
only one of the three that actually attempts and roughly recovers something close to the real 6-way
activity structure (ARI 0.418, NMI 0.563), even if it's far from perfect. DBSCAN's headline internal
scores are misleading once you check them against ARI/NMI -- it settled for a coarse 2-cluster
solution that trades away the fine-grained structure entirely.
\item \textbf{How sensitive was K-Means to the choice of k:} Very. The silhouette score drops sharply
from 0.4775 at k=2 to 0.1835 at k=6, meaning the natural, cleanly-separated structure in this data is
really just 2 clusters (moving vs stationary); forcing k=6 finds real sub-structure but at a real cost
to how cleanly separated the individual clusters are.
\item \textbf{Did DBSCAN detect noise or small clusters effectively:} It found noise, plenty of it,
but not small clusters -- every $\epsilon$ that avoided lumping everything into one cluster still only
produced 2 clusters, never anything close to 6. It's better described here as a noise/outlier
detector than as a way to recover the six activities.
\item \textbf{How does linkage choice affect Hierarchical Clustering:} Hugely. Single linkage chains
points together and basically produces one dominant cluster plus a scattering of near-singleton
outlier clusters, which inflates the silhouette score without meaning much. Complete and average
linkage sit in between. Ward's linkage is the only one that produces genuinely balanced, compact
clusters at k=6, which is presumably why the experiment specifically calls for it, even though its raw
silhouette number is the lowest of the four.
\item \textbf{Which internal metric best matched visual intuition:} Calinski-Harabasz tracked the
t-SNE/PCA plots better than silhouette did in this case -- K-Means' relatively high CH score (858.3)
lines up with clusters that look reasonably compact and separated in the scatter plots, whereas
silhouette got fooled by DBSCAN's 2-cluster solution looking artificially clean.
\end{itemize}

\section{Discussion}
The big theme here is that "looks good by an internal metric" and "actually recovers the known
structure" are not the same thing, and DBSCAN is the clearest example: best silhouette and
Davies-Bouldin scores of the three, but worst ARI, because it quietly gave up on separating six
activities and settled for two. K-Means, despite having the worst silhouette score of the three at
k=6, is the one that comes closest to matching the ground truth, simply because it was forced to
actually produce six clusters rather than finding an easy two-way split. The SITTING/STANDING overlap
that trips up every algorithm here isn't really a clustering failure so much as a property of the data
itself -- those two activities produce genuinely similar sensor signatures when a phone is worn at the
waist, so no amount of retuning $\epsilon$ or trying a different linkage is likely to fully separate
them without extra features or a different sensor placement.

\section{Conclusion}
Across K-Means, DBSCAN and Hierarchical (Ward) clustering on the HAR dataset, K-Means at $k=6$ gave the
most useful match to the real six activities (ARI 0.418, NMI 0.563), mainly by being willing to accept
a lower silhouette score in exchange for actually splitting the data six ways. DBSCAN is genuinely good
at finding the coarsest possible structure (moving vs. stationary) and flagging outliers, but it isn't
well suited to recovering finer-grained classes on this dataset. Ward's linkage HAC lands in between
the two, doing reasonably on external metrics without either of K-Means' flexibility or DBSCAN's noise
handling. If the actual goal is to recover activity-level clusters, K-Means with a task-informed choice
of $k$ (rather than the value silhouette alone would suggest) is the more practical choice here.

\section{Learning Outcomes}
\begin{itemize}
\item Learned to read an elbow curve and a silhouette-vs-k plot together, and to recognize when they
disagree with what the actual task needs (silhouette wanted $k=2$, the task wanted $k=6$).
\item Got hands-on practice tuning DBSCAN's $\epsilon$ using a k-distance plot, and saw first-hand how
narrow the useful range of $\epsilon$ can be on high-dimensional, PCA-reduced sensor data.
\item Directly observed the single-linkage chaining problem in hierarchical clustering, and understood
why Ward's linkage is usually preferred for producing balanced clusters even when its raw silhouette
looks worse.
\item Practiced distinguishing internal clustering metrics (silhouette, Davies-Bouldin,
Calinski-Harabasz) from external ones (ARI, NMI), and saw a concrete case where they disagreed.
\item Used PCA and t-SNE together for visualization, and noticed how t-SNE separates the
SITTING/STANDING overlap a little better than a straight 2D PCA projection does.
\end{itemize}

\section{References}
\begin{enumerate}
\item Davide Anguita, Alessandro Ghio, Luca Oneto, Xavier Parra and Jorge L. Reyes-Ortiz, ``Human
Activity Recognition on Smartphones using a Multiclass Hardware-Friendly Support Vector Machine'',
IWAAL 2012.
\item UCI Machine Learning Repository, Human Activity Recognition Using Smartphones Data Set:
\url{https://archive.ics.uci.edu/dataset/240/human+activity+recognition+using+smartphones}
\item Scikit-learn documentation: \url{https://scikit-learn.org/stable/modules/clustering.html}
\end{enumerate}

\end{document}
